% A short, plain-language explainer of the changes recently made to
% Theseus. Compiled via:
%   pdflatex docs/guides/07_Theseus_After_Round19.tex
% (run twice for cross-references).

\documentclass[11pt]{article}
% Shared preamble for Theseus user guides.
% Compiled with pdflatex. See docs/guides/BUILD.md for rationale.

\usepackage[a4paper,margin=0.85in]{geometry}
\usepackage[T1]{fontenc}
\usepackage[utf8]{inputenc}

% Body: Palatino (via mathpazo). Headings: Helvetica (sans).
% Palatino is requested in the house style as a substitute for Charter
% because Charter is not part of the standard TeX Live distribution
% that the build container ships.
\usepackage{mathpazo}
\usepackage[scaled=0.92]{helvet}
\usepackage{courier}
\renewcommand{\familydefault}{\rmdefault}

\usepackage{microtype}
\usepackage{parskip}
\usepackage{enumitem}
\usepackage{booktabs}
\usepackage{longtable}
\usepackage{titlesec}
\usepackage{xcolor}
\usepackage{fancyhdr}
\usepackage{fancyvrb}
\usepackage{graphicx}
% Allow `_` in text mode without escaping. The screenshot file names and
% some environment-variable references in the guides contain underscores.
\usepackage{underscore}
\usepackage{tcolorbox}
\tcbuselibrary{breakable,skins}
\usepackage[hidelinks,colorlinks=true,linkcolor=black,urlcolor=teal,citecolor=black]{hyperref}

% Sans for chapter and section headings; serif body keeps prose readable
% on screen and in print.
\titleformat{\section}{\sffamily\Large\bfseries}{\thesection}{0.7em}{}
\titleformat{\subsection}{\sffamily\large\bfseries}{\thesubsection}{0.6em}{}
\titleformat{\subsubsection}{\sffamily\normalsize\bfseries}{\thesubsubsection}{0.5em}{}

\pagestyle{fancy}
\fancyhf{}
\fancyhead[L]{\small\sffamily Theseus User Guide}
\fancyhead[R]{\small\sffamily\thepage}
\renewcommand{\headrulewidth}{0.3pt}

\setlength{\parindent}{0pt}

% Convenience commands shared by every guide.
\newcommand{\page}[1]{\textbf{#1}}
\newcommand{\file}[1]{\texttt{#1}}
\newcommand{\route}[1]{\texttt{#1}}
\newcommand{\env}[1]{\texttt{#1}}
\newcommand{\code}[1]{\texttt{#1}}

% Callout boxes used by every guide.
\newtcolorbox{tldr}{
  colback=teal!5!white,colframe=teal!55!black,
  fonttitle=\sffamily\bfseries,title=TL;DR,
  breakable,boxrule=0.4pt,arc=2pt
}
\newtcolorbox{youwillneed}{
  colback=gray!4!white,colframe=gray!55!black,
  fonttitle=\sffamily\bfseries,title=You will need,
  breakable,boxrule=0.4pt,arc=2pt
}
\newtcolorbox{notcovered}{
  colback=orange!5!white,colframe=orange!55!black,
  fonttitle=\sffamily\bfseries,title=This guide does NOT cover,
  breakable,boxrule=0.4pt,arc=2pt
}
\newtcolorbox{preview}{
  colback=yellow!8!white,colframe=yellow!55!black,
  fonttitle=\sffamily\bfseries,title=Preview --- partially shipped,
  breakable,boxrule=0.4pt,arc=2pt
}
\newtcolorbox{nextup}{
  colback=blue!4!white,colframe=blue!55!black,
  fonttitle=\sffamily\bfseries,title=Where to read next,
  breakable,boxrule=0.4pt,arc=2pt
}

% Insert a screenshot only when the file exists, so the build does not
% break before the Playwright capture run is wired up.
\newcommand{\screenshot}[3]{%
  \IfFileExists{screenshots/#1}{%
    \begin{figure}[h]\centering
      \includegraphics[width=\linewidth]{screenshots/#1}
      \caption{#2}\label{fig:#3}
    \end{figure}%
  }{%
    \begin{center}\small\itshape%
      [screenshot \texttt{screenshots/#1} not yet captured --- #2]%
    \end{center}%
  }%
}


\title{\sffamily\textbf{What Changed in Theseus} \\[0.4em]
       \large\sffamily A Plain-Language Explainer}
\author{}
\date{2026-05-16}

\begin{document}
\maketitle
\thispagestyle{empty}

\begin{tldr}
Theseus was rebuilt around a new idea: that what the system stores
should be \emph{principles} (generalizable truth claims) rather than
\emph{summaries} (what the author of a document seems to think). On
top of that change, a few related pieces were added — a way to
turn principles into structured reasoning, a cleaner contradiction
engine, a way to mark each source by where it came from, an
output format for the firm's takes, and a few read views. This
explainer says what each change is and why it was made — nothing
about the underlying infrastructure. The final section explains
how to re-process the existing corpus, which is the one thing
that requires your action.
\end{tldr}

\tableofcontents
\newpage

% ---------------------------------------------------------------------------
\section{The core problem we fixed}
% ---------------------------------------------------------------------------

The old system read a document and stored what it thought the
author was saying \emph{about themselves} — sentences like ``I
have since become a fan of Peter Thiel's idea.'' That kind of
conclusion is autobiographical. The firm can't act on it. We can
only act on what an idea \emph{is} and \emph{when it applies}.

So the entire pipeline was reshaped around principles: third-person,
generalizable claims about how the world works. ``Definite optimism
outperforms indefinite optimism over a $\geq 10$-year horizon,
conditional on a coherent vision'' is a principle. The firm
\emph{can} act on that — it can identify situations that match the
condition and form positions about what's likely to follow.

Everything else in this guide either follows from that change or
cleans up something that was getting in the way of it.

% ---------------------------------------------------------------------------
\section{What changed}
% ---------------------------------------------------------------------------

\subsection{Conclusions are now principles}

The extractor no longer emits first-person summaries. If a source
span is purely autobiographical and no underlying principle can be
extracted from it, the extractor logs that fact and emits nothing.
If a principle is present, it's stored with a few extra pieces of
structure: what \emph{kind} of principle it is (rule, criterion,
mechanism, heuristic, definition, formula, algorithm), what domain
it applies to, what measurable quantities could test it, and a few
concrete decisions it would inform.

This is the load-bearing change. The existing corpus needs to be
re-processed through the new extractor — see
Section~\ref{sec:reprocessing} below.

\subsection{A new layer on top of principles: algorithms}

A principle says \emph{what is true}. An algorithm says \emph{when,
given what we can observe, the principle predicts something
specific}. So an algorithm is the bridge between abstract principle
and concrete prediction.

Each algorithm names the observable things it watches, a condition
that has to be true for it to fire, the principles it's reasoning
from, and the structured output it produces. The reasoning chain
inside an algorithm cites the principles it depends on, so you can
always trace a prediction back to its sources.

A canonical example: an ``arms-race escalation'' algorithm watches
bilateral military-spending changes, rhetoric levels, and whether
a third-party mediator is present. When the conditions match, it
applies two principles about security dilemmas and emits a
predicted spending increase with confidence bands.

Algorithms are drafted from clusters of related principles. You
review each draft, accept/edit/reject it, and from then on the
algorithm runs whenever its conditions are met.

\subsection{One contradiction engine, not six}

The old system had six different ways of detecting that two
claims contradict each other — different heuristics that
disagreed with each other and were hard to calibrate. That was
replaced with a single detector that returns a calibrated score,
a confidence band, and a human-readable explanation of where the
disagreement lies.

A second piece of this change: contradictions no longer have a
``resolve'' button. Operator clicks were treated as authoritative,
which was wrong — the database should reflect what the sources
jointly imply, not which side an operator preferred on a given
afternoon. Contradictions now resolve when new source material
weights one side decisively over the other.

\subsection{Provenance is set at upload time}

Every piece of source material now carries one of four labels,
chosen when it's uploaded:

\begin{itemize}[leftmargin=*]
  \item \textbf{Proprietary} — the firm wrote it.
  \item \textbf{Endorsed external} — someone else wrote it, but we
        explicitly endorse it as representative of our thinking.
  \item \textbf{Studied external} — reference material; we read it
        but don't endorse it.
  \item \textbf{Opposing external} — material we explicitly
        disagree with, kept for the value of testing our positions
        against it.
\end{itemize}

Why this matters: the system used to flag contradictions between
the firm's principles and opposing material it was reading — which
is noise, because we \emph{expect} to disagree with that material.
Provenance demarcation lets the contradiction engine skip those
cross-provenance pairs and surface only the contradictions that
are actually news.

\subsection{Memos as the canonical output}

When the system answers a question — whether it's ``what's our
take on X?'' or ``should we bet on Y?'' — the answer now takes a
fixed memo shape: a TL;DR, the question, the governing principles,
the observed inputs, the reasoning chain, the implied position,
what would change our mind, and any caveats or abstentions. The
memo is the firm's canonical output.

The system explicitly prefers to \emph{abstain} over making a
chain of reasoning it can't ground. If there are no governing
principles, or the principles directly contradict each other, or
the confidence band is too wide, or the question itself is
unformed — the memo says so, by name.

\subsection{``Bet'' was generalized}

A bet used to mean a financial position on a prediction market or
an equity. ``Bet'' now means any falsifiable commitment of firm
resources: a financial position on a market, a public statement of
position with no money behind it, an internal allocation of
operator time or hiring direction, or a scientific prediction that
resolves against external data. All four are tracked in one place.

The reason: the firm's edge is the principle layer, not any
specific way of expressing a view. Restricting bets to financial
markets undersold what the system could do.

\subsection{A knowledge graph view}

The corpus is no longer a flat list of claims. There's now a
cross-source graph view where principles, sources, algorithms,
memos, and concepts are nodes, and the relationships between them
are edges (``derived from'', ``contradicts'', ``supports'',
``applies to'', ``cites''). Click any edge and a reasoner produces
a grounded explanation of why that edge exists. The reasoner
refuses to fabricate connections the data doesn't support.

This is the reading view for ``what does the corpus jointly imply
about X?'' — a question that wasn't directly answerable before,
because the underlying structure didn't exist.

\subsection{Live conversation recording}

The meeting-recording surface (``Dialectic'') now fires real-time
contradiction flags as a participant says something that
contradicts (a) something they or another participant said earlier
in the same conversation, or (b) a principle the firm has already
committed to. Provisional principles can be extracted live and
triaged after the session ends.

The motivation: the most valuable corpus material is often the
live conversation, but it was historically lost or transcribed
only as an afterthought. Real-time flags turn the conversation
itself into a feedback loop with the firm's existing commitments.

\subsection{Surface cleanup}

A few smaller changes worth knowing about: the homepage and
about page now lead with the ``Philosopher in a Box'' positioning
— Theseus is the firm's machine, not a product for sale. The old
``Attention'' panel on the dashboard is gone (redundant once the
Currents feed and algorithm-invocation feeds matured). The legacy
contradiction heuristics are marked deprecated and no longer write
new rows.

% ---------------------------------------------------------------------------
\section{Re-processing the existing corpus}
\label{sec:reprocessing}
% ---------------------------------------------------------------------------

\begin{tldr}
This is the one thing that requires your action. Existing
conclusions are NOT auto-replaced — the new extractor's output
goes to a triage queue at \route{/(authed)/extractor/re-extract}
where you accept, edit, or reject each replacement. Do a small
sample first to confirm the new extractor behaves as you want.
\end{tldr}

\subsection{Step 1: sanity-check the new extractor}

Before queueing the whole corpus, run the new extractor against
one or two artifacts whose existing conclusions had the
first-person symptom (e.g., the one that produced ``I have since
become a fan of Peter Thiel's idea''). The single-artifact dry-run
prints the proposed principles without persisting anything:

\begin{verbatim}
cd ~/Desktop/Theseus
source .venv-currents/bin/activate
python noosphere/scripts/reextract_corpus.py \
       --artifact <artifact_id> --dry-run
\end{verbatim}

Confirm the output: each proposed principle should be third-person
and generalizable, the principle-kind should match your reading,
the source citation should be intact, and refusals should fire on
spans that are purely autobiographical.

\subsection{Step 2: dry-run a small sample}

Once one or two artifacts look right, sample five at random to
confirm the extractor is consistent:

\begin{verbatim}
python noosphere/scripts/reextract_corpus.py --sample 5 --dry-run
\end{verbatim}

If any of the five looks wrong, stop and adjust the principle
extractor's system prompt before going further.

\subsection{Step 3: queue the whole corpus}

When the sample looks right, queue every artifact for re-extraction.
You can do this all at once, or restrict by date or by provenance:

\begin{verbatim}
# Everything:
python noosphere/scripts/reextract_corpus.py --all

# Only your own writing first (highest signal, lowest noise):
python noosphere/scripts/reextract_corpus.py --all \
       --provenance PROPRIETARY

# Only material added since 2026-03-01:
python noosphere/scripts/reextract_corpus.py --all \
       --since 2026-03-01
\end{verbatim}

The job is asynchronous and respects the system's hourly budget,
so a large corpus may take several cycles to drain.

\subsection{Step 4: triage the queue}

Each proposed replacement lands at
\route{/(authed)/extractor/re-extract}, NOT as an automatic
overwrite. You review each row, which shows:

\begin{itemize}[leftmargin=*]
  \item The original source span.
  \item The existing first-person conclusion (from the old extractor).
  \item The new principle-shaped proposal.
  \item Three buttons: \textbf{Accept}, \textbf{Edit \& Accept}, or
        \textbf{Reject}.
\end{itemize}

The system never replaces anything on its own — every row needs
your confirmation. Once you accept a row, the new principle gets
embedded, joins a cluster, gets tested for contradictions against
its siblings, and becomes searchable on the principles page.
Algorithms may be drafted from new clusters automatically.

\subsection{What it'll cost, roughly}

For a corpus of $\sim$500 artifacts:

\begin{itemize}[leftmargin=*]
  \item \textbf{Tokens.} Roughly 1.5--2M total tokens.
  \item \textbf{Wall clock.} At the default hourly budget, expect
        6--24 hours for the queue to drain (smaller corpora
        proportional).
  \item \textbf{Your triage time.} Each row takes 30--60 seconds
        if the proposed principle is reasonable. For 500 artifacts,
        budget 4--8 hours of triage, spread across whichever days
        work. There's no rush.
\end{itemize}

% ---------------------------------------------------------------------------
\section{Where to read next, if you want depth}
% ---------------------------------------------------------------------------

The earlier guides in this series go deeper on individual surfaces:

\begin{itemize}[leftmargin=*]
  \item \file{docs/guides/01_Theseus_Quick_Start.pdf} — the minimal
        intro for a new reader.
  \item \file{docs/guides/02_Knowledge_and_Principles.pdf} — how
        ingestion-to-principles works.
  \item \file{docs/guides/03_The_Oracle.pdf} — using the Oracle.
  \item \file{docs/guides/04_Currents.pdf} — the live opinion feed.
  \item \file{docs/guides/05_Forecasts_and_Portfolio.pdf} —
        prediction markets and equities.
  \item \file{docs/guides/06_Operator_Console.pdf} — the operator
        surface and the safety architecture in detail.
\end{itemize}

\end{document}
